%% kapa: exam=1 round=28 year=2017 period=2 form=A subject=accounting subject_ko=회계학 seq=7 number=47
%% source: Q-Net 감정평가사 기출문제 (cst00309) artlSeq=5199603 file=extracted/1008/5199603/2017년 제28회 감정평가사 1차시험 문제지/2교시 A형.pdf
%% conversion: kordoc 4.24.51 -> tools/gen1.py (auto); review=proofread-sonnet; flags=none
\begin{question}{47}
20x1년 초 (주)감평은 정부보조금 ￦500,000을 받아 연구소 건물(내용연수 5년, 잔존가치 ￦0, 정액법 상각)을 ￦1,000,000에 취득하고 다음과 같이 회계처리를 하였다.
\begin{framedtext}
\centering (차변) 건물 1,000,000\qquad\qquad (대변) 현금 1,000,000
\end{framedtext}
위 거래와 관련하여 정부보조금 및 감가상각에 대한 회계처리가 누락되었다. 이를 장부마감 이전에 반영하여 재무제표에 표시할 경우, 20x1년 말 재무제표에 미치는 영향으로 옳은 것은? (단, 정부보조금은 자산에서 차감하는 방법으로 회계처리한다.) \par
\begin{choices}
  \item 20x1년말 총자산금액에는 영향이 없음
  \item 20x1년도 당기순이익은 ￦100,000 감소
  \item 20x1년말 총부채 금액은 ￦400,000 증가
  \item 20x1년말 감가상각누계액은 ￦200,000 감소
  \item 20x1년말 건물 장부금액은 ￦200,000 감소
\end{choices}
\end{question}
\answer{2}
